\section{Galois slicing and automatic differentiation via categories of families}
\label{sec:galois-slicing-auto-diff-via-fam}

Recall that $\Fam(\cat{C})$ is the Grothendieck construction $\Grothendieck{X}\Fam(X, \cat{C})$.

\subsection{\GPS}
\label{sec:galois-slicing-auto-diff-via-fam:galois-slicing}

$\Fam(\LatGal)$ has:
\begin{itemize}
\item as objects $(X, \partial X)$ all pairs of a set $X$ and and for every $x \in X$, a bounded lattice
$\partial X_x$;
\item as morphisms $(X, \partial X) \to (Y, \partial Y)$, all pairs $(f, \partial f)$ of a function $f: X \to
Y$ and for every $x \in X$, a Galois connection $\partial f_x: \partial X_x \to \partial Y_{f(x)}$.
\end{itemize}

\noindent with the role of the pullback along $f$ in the composition $(g, \partial g) \comp (f, \partial f) =
(g \comp f, \reindex{\partial g}{f} \comp \partial f)$ being to select the appropriate lattice of
approximations and Galois connection at each point.

\subsection{Automatic differentiation}
\label{sec:galois-slicing-auto-diff-via-fam:auto-diff}

Consider the functor $\Diff \hookrightarrow \Fam(\FinVect)$ which send objects $\RR^n$ to constant families of
tangent spaces $(\RR^n, \const(\RR^n))$ and morphisms $f: \RR^m \to \RR^n \times (\RR^m \linearto \RR^n)$ to
pairs $(f_1, f_2)$ of differentiable functions $f_1: \RR^m \to \RR^n$ and families of linear maps $f_2:
\{{f_2}_x: \RR^m \linearto \RR^n\}_{x \in \RR^m}$. This exhibits $\Diff$ as equivalent to a full subcategory
of $\Fam(\FinVect)$.

In that subcategory, for any $(f, \partial f): (\RR^m, \const(\RR^m)) \to (\RR^n, \const(\RR^n))$ and $(g,
\partial g): (\RR^n, \const(\RR^n)) \to (\RR^k, \const(\RR^k))$, we have that for any $x \in \RR^m$:
\begin{itemize}
\item $\partial f_x: \RR^m \to \RR^n$;
\item $\reindex{\partial g}{f}_x = \partial g_{f(x)}: \RR^n \to \RR^k$.
\end{itemize}

\noindent The composition $(g, \partial g) \comp (f, \partial f) = (g \comp f, \reindex{\partial g}{f} \comp
\partial f): (\RR^m, \const(\RR^m)) \to (\RR^k, \const(\RR^k))$ corresponds exactly to the forward chain rule,
since $(\reindex{\partial g}{f} \comp \partial f)_x = {\partial g}_{f(x)} \comp {\partial f}_x$.

Similarly, $\Diff^*$ is equivalent to a full subcategory of $\Fam(\FinVect^\op)$ with composition
corresponding to the backward chain rule.

\subsection{Rational intervals and two forward analyses}
\label{sec:galois-slicing-auto-diff-via-fam:rational-intervals}

A concrete instance illustrating the contrast between the Galois (right adjoint) forward map and the Tarski
conjugate forward map is given by intervals over the rationals. For each $q \in \QQ$ let $\Intv(q) = \{[l, u] :
l \leq q \leq u\}$ ordered by reverse inclusion (so a tighter interval is \emph{higher} in the order, i.e.\ the
information order). Lifting by an explicit bottom yields a bounded distributive lattice $\IntvL(q) =
\Intv(q)_\bot$, and the family $q \mapsto \IntvL(q)$ is an object of both $\Fam(\LatGal)$ and $\Fam(\LatConj)$.

Addition has three associated maps for any $q_1, q_2$:
\begin{itemize}
\item $\add: \IntvL(q_1 + q_2) \to \IntvL(q_1) \times \IntvL(q_2)$ — the
  \emph{join-preserving backward} map, defined by subtraction: $\add(z).\proj_1 = [z.l - q_2, z.u - q_2]$. This
  is shared between Galois and conjugate (left of both adjoints).
\item $\add^*: \IntvL(q_1) \times \IntvL(q_2) \to \IntvL(q_1 + q_2)$ — the \emph{meet-preserving forward}
  (Galois right adjoint), with $\add^*(x, y).l = (q_2 + x.l) \meet (q_1 + y.l)$, $\add^*(x, y).u = (q_2 + x.u)
  \join (q_1 + y.u)$.
\item $\addT: \IntvL(q_1) \times \IntvL(q_2) \to \IntvL(q_1 + q_2)$ — the \emph{join-preserving forward}
  (Tarski conjugate of $\add$), with $\addT(x, y).l = (q_2 + x.l) \join (q_1 + y.l)$, $\addT(x, y).u = (q_2 +
  x.u) \meet (q_1 + y.u)$.
\end{itemize}

\noindent Each input interval $x \in \Intv(q_i)$ encodes a constraint of the form ``the $i$-th summand lies in
$[x.l, x.u]$''; the maps $\add^*$ and $\addT$ shift each input by the \emph{other} summand's index to obtain
candidate ranges for the result, then combine them.

\paragraph{Comparison.} On overlapping but non-nested input intervals (e.g. $\IntvL$-shifted ranges $[1/2, 1]$
and $[4/5, 3/2]$):
\[
  \addT\colon\quad \big[\max(1/2, 4/5),\; \min(1, 3/2)\big] = [4/5, 1] \qquad \text{(intersection)}
\]
\[
  \add^*\colon\quad \big[\min(1/2, 4/5),\; \max(1, 3/2)\big] = [1/2, 3/2] \qquad \text{(union)}
\]

\noindent So $\addT$ \emph{tightens}: it returns the smallest interval simultaneously consistent with all
inputs (the join in the information order), modelling ``each input is evidence; combine evidence by joining
constraints''. $\add^*$ \emph{broadens}: it returns the smallest interval covering all inputs (the meet),
modelling ``each input is one possibility; combine possibilities by overapproximation''.

\paragraph{Connection to AD.} Forward-mode AD propagates tangents through addition by $dy = dx_1 + dx_2$.
Modulo idempotence of $+$ (replacing it by $\join$ in the information order), this corresponds to $\addT$:
each input contributes a constraint on the output, and the constraints are joined into the strongest
common refinement. By contrast, $\add^*$ is the abstract-interpretation flavour, propagating sets of possible
values via overapproximation.
